% Dot batch 85; source SHA-256: 1304dece8a13b8c523e173f4655af74106bd2635d4e0f9f199b19cda7e0e8f99
% Archived source: dot_updates_20261124/batches/batch085/EGA第五卷_线性系的除子集合与底除子.tex
% Original body lines: 27--75; integration map: integration/late_integration.json
\begingroup
\def\newscope{V-L}
\def\Div{\operatorname{div}}
\def\Ens{\operatorname{Ens}}
\def\OO{\mathcal O}
\def\KK{\mathcal K}
\def\PP{\mathbf P}
再回到地面上，甚至再贴近地面一些：为了用通俗的话说明线性系的概念，
姑且取一个基域以求简便（尽管换成仿射底，陈述实质上仍可保持原样）。\cmt{L18-085-01}{法146的基域为 \emph{corps de base}；英110的“box field”非数学术语。关于仿射底的话是作者的提示，不是此处已经证出的推广。}

\egas{4.4}{命题}

设 \(X\) 是域 \(k\) 上有限型的预概形，且自然映射
\[
 k\longrightarrow H^0(X,\OO_X)
\]
是同构。对 \(X\) 上由 \(k\) 上射影丛 \(P^\vee\) 参数化的每个除子线性系 \(D\)，
记 \(\Ens(D)\) 为形如 \(D(t)\) 的除子组成的集合（注意：是集合！），其中
\(t\in P^\vee(k)\)。\cmt{L18-085-02}{\(\Ens(D)\) 只取 \(k\)-有理参数，归入 \(X\) 上的 Cartier 除子集合；不是取线性等价类，也不是命题4.3中的除子子函子。若允许零维参数向量空间，约定其射影空间为空；若只定义正秩射影丛，则以下线性系的除子集合非空。}

\smallskip\noindent\textbf{a)}
设 \(D'\) 是 \(X\) 上另一个除子线性系，由射影丛 \((P')^\vee\) 参数化。
则 \(D\) 与 \(D'\) 同构，当且仅当 \(\Ens(D)=\Ens(D')\)。

\smallskip\noindent\textbf{b)（校订）。}
假设 \(k\) 代数闭，或 \(X\) 几何整。\cmt{L18-085-03}{法146边注补入此二选一条件；b) 随后的“为空，或……”被划去，并改接法147的向量子空间条件。此处不将删文复作正文。英111保留了边注的两项假设，但省略了若干手稿改写痕迹。}
记 \(\KK_X\) 为全分式层。\(X\) 上非空的有效 Cartier 除子集合 \(\Delta\) 形如
\(\Ens(D)\)，当且仅当存在有效 Cartier 除子 \(A\) 和有限维非零
\(k\)-向量子空间 \(E\subseteq\Gamma(X,\KK_X)\)，使每个 \(f\in E\setminus\{0\}\)
都属于 \(\Gamma(X,\KK_X^\times)\)，且
\[
 \Delta=\{\, A+\Div(f): f\in E\setminus\{0\}\,\}.
\]
右侧中的每个除子必须有效。可以取 \(A\in\Delta\)；以它为底时，可以取 \(1\in E\)。
这里 \(A\) 只是用来比较的基准除子，不必是线性系共有的固定部分；\(\Div(f)\) 是主 Cartier 除子，允许有极点。\cmt{L18-085-04}{法147改称亚纯函数，\emph{régulière} 由“\(\Div(f)\) 有定义”解释，即在全分式层中可逆；不是说函数处处无极点。手稿与英111都把 \(\Delta\) 直接写成 \(\{\Div(f)\}\)，未写平移项，也未明确限制 \(E\) 有限维。按字面，若 \(\Div(f)\) 有效，\(f\) 局部是非零因子与单位的乘积，故 \(f\in H^0(X,\OO_X)=k\)，只能给零除子。反例：\(\PP^1_k\) 上 \(\OO(1)\) 的完全线性系包含一次点除子。正文采用可由4.1严格推出的修订，不把它冒称手稿原句。}

\smallskip\noindent\textbf{c)（校订）。}
设 \((A,E)\)、\((A',E')\) 满足 b) 的条件，并分别给出 \(\Delta\)、\(\Delta'\)。
则 \(\Delta=\Delta'\)，当且仅当存在 \(h\in\Gamma(X,\KK_X^\times)\)，使
\[
 E'=hE,\qquad A'=A-\Div(h).
\]
存在时，\(h\) 在乘以 \(k^\times\) 的元素的意义下唯一。\cmt{L18-085-05}{法147的 c) 仍印有“严格有理函数”，对乘子又改作“正则亚纯函数”；英111相应使用 \emph{regular pseudo-function}，并误排 \(k^\times\)。手稿只写乘法条件 \(E'=hE\)，未写除子的抵消关系。若严格要求两组主除子都有效，则 b) 的字面表述只剩平凡情形。若放宽为一般主除子来理解 c) 的乘法断言，取 \(X=\PP^1_k\)、\(E=k\)、\(E'=kx\)、\(h=x\)，则两个集合分别是 \(\{0\}\) 与 \(\{[0]-[\infty]\}\)。反过来，不带平移的主除子集合若相等，则每对同除子函数之比是全局单位，在本命题下属 \(k^\times\)，所以 \(E=E'\)。}

\medskip\noindent\textbf{作者说明（译）。}
利用4.1，证明就是一个简单的练习；除非你要求，我就不写出来了。
此外我觉得，把4.4移到4.3前面也许更好，因为在技术上它要简单得多。
还应注意，若 \(X\) 几何整，b) 对 \(E\) 所提的那项条件自动满足。\cmt{L18-085-06}{指非零亚纯函数为全分式层的单位这一条件；对于校订后的表述，
有限维性以及 \(A+\Div(f)\) 的有效性仍须满足。}

在 b) 开头加限制，是因为否则，b) 所述条件在过到 \(k\) 的代数闭包后可能不再成立。
在任一特征中都可以举例，甚至 \(k\) 在正特征下可分闭时也这样。
要把话说周全，b) 应不附加关于 \(k\) 或 \(X\) 的条件，而是在过到 \(k\) 的代数闭包后
陈述对 \(E\) 的条件（并指出在几何整情形此项条件自动满足）。
稍微滥用一下术语，形如 \(\Ens(D)\) 的除子集合 \(\Delta\)，常径称为 \(X\) 上的除子线性系。\cmt{L18-085-07}{最后一句为法147手写；英111标作略难辨读，并在其后加了手稿正文告终的编辑说明和拉丁文结束语。法148另起书信。这里的“在代数闭包上检验”是检验基变换后的向量空间及其全部新线性组合，不是只检查原有 \(k\)-有理参数。}
\par\endgroup
